\documentclass[../main.tex]{subfiles}
\begin{document}
%%%%%%%%%%%%%%%%%%%%%%%%%%%%%%%%%%%%%%%%%%%%%%%%%%%%%%%%%%%%%%%%%%%%%%%%%%%%%%%%%%%%%%%%%%%%%%%%%%%
\chapter{Diagramming Basics}

Now that you understand the basics of working with OmniGraffle, it's time for a quick tutorial.
The purpose of this tutorial is to help you master OmniGraffle as soon as possible.

OmniGraffle excels at clarifying complex relationships, and nothing says ``complex
relationships'' like a Shakespeare romantic comedy. Let's try to make some sense of
\textit{Much Ado About Nothing}, and learn a bit about OmniGraffle in the process.

\ogfigure{0.800}{res/og7_diagramming51.png}

Of course, this play is just a handy example; if there's another story you're more familiar
with, you can go ahead and use that instead.

So grab your favorite beverage, find something soothing to listen to, and let's get to work!

\section{Setting Up Your OmniGraffle Project}

Start by creating a new OmniGraffle document:

\begin{enumerate}
  \item Choose \textbf{File} $\blacktriangleright$ \textbf{New} (\textbf{Command-N}). This
    creates a new OmniGraffle document using the default template. Your OmniGraffle window
    should look like this:

    \ogfigure{0.800}{res/og7_diagramming01.png}

    A nice, blank slate to work on!
  \item Next, press \textbf{Command-3} to switch the Inspector sidebar on the right to the
    \textbf{Canvas} inspectors.
  \item \textbf{Option-click} on the \textbf{Diagram Layout and Style} inspector; this opens the
    Diagram Layout and Style inspector, while collapsing the others in the Inspector sidebar.

    For the purpose of this tutorial, leave the diagram style set to \textbf{Hierarchical}, but go
    ahead and turn on \textbf{Auto layout} by clicking the checkbox, as so:

    \ogfigure{0.800}{res/og7_diagramming05.png}

    Auto layout ensures that the objects in the diagram you create are evenly spaced, and that
    they automatically flow into place if you change the diagram style or add new items to the
    diagram.
\end{enumerate}

\begin{ogtip}
  You can learn more about the Inspectors in Working With the Inspectors, later in Part II.
\end{ogtip}

With the project set up, choose \textbf{File} $\blacktriangleright$ \textbf{Save} to save the
OmniGraffle file somewhere on your Mac. Be sure to save your work at various points along the
way, or before moving on to the next step in the tutorial.

\section{Defining a Base Shape for Your Diagram}

The next bit of setup that you'll want to do is define a shape to use as the default for the
diagram.

\begin{enumerate}
  \item Select the \textbf{Shape} tool in the Tool Palette:

    \ogfigure{0.800}{res/og7_diagramming06.png}
  \item Press \textbf{Command-1} to switch to the \textbf{Object} inspectors.
  \item \textbf{Option-click} the \textbf{Shape} inspector.
  \item Click the \textbf{Shape Selector} button at the top of the Shape inspector. This opens a
    popover menu from which you can choose from a variety of default shapes included with
    OmniGraffle.
  \item In the upper-left corner, double-click the \textbf{Rectangle} shape:

    \ogfigure{0.800}{res/og7_diagramming07.png}
  \item To the left of the Tool Palette, click the \textbf{Style Well}. This opens a pop-up menu
    with details about the shape to be drawn; in this case, a rectangle with a white fill and a
    a simple black stroke as its outline:

    \ogfigure{0.800}{res/og7_diagramming08.png}
\end{enumerate}

\section{Creating Objects in the Sidebar's Outline Tab}

To create and instantly label the objects for the \textit{Much Ado About Nothing} diagram,
you'll use the \textbf{Outline} tab in the Sidebar.

\begin{enumerate}
  \item Press \textbf{Option-Command--4} to switch the Sidebar to the \textbf{Outline} tab.
  \item Click where it says \textbf{Click to add a topic} to create your first object. A
    \textbf{Topic} item appears in the outline, and an object labeled \textbf{Topic} appears on
    the canvas simultaneously.

    \ogfigure{0.800}{res/og7_diagramming10.png}
  \item The item's text is already highlighted, so you can just start typing to rename it; enter
    \textbf{Don John}, and then press \textbf{Return} to finish with this item and move on to
    the next one.

    \ogfigure{0.673}{res/og7_diagramming11.png}
  \item A new item appears below the first one in the Outline tab, and a new shape appears on
    the canvas to the right; name this object \textbf{Don Pedro} and press \textbf{Return}.
  \item Name the third object \textbf{Leonato} and then press \textbf{Escape}:

    \ogfigure{0.800}{res/og7_diagramming12.png}

    Unlike pressing \textbf{Return}, which creates another object in the outline, pressing
    \textbf{Escape} accepts the text you entered. (You can also press \textbf{Command-Return} to
    turn editing on and off for the selected object.)
\end{enumerate}

Now you have three items in your Outline and three shapes on your canvas. Time to add some more
structure to the diagram!

\section{Creating Depth by Adding Child Objects}

All of the characters you've created so far exist on the same level, but now let's add some more
characters one level below them. The next two characters are soldiers who fought under Don
Pedro's command.

\begin{enumerate}
  \item Click to select \textbf{Don Pedro} in the Outline tab, and press \textbf{Return} to
    create a new object.
  \item Press \textbf{Tab} to indent this item underneath Don Pedro, and name it
    \textbf{Benedick}:

    \ogfigure{0.800}{res/og7_diagramming13.png}
  \item Press \textbf{Return} to create another item at the same level, name it
    \textbf{Claudio}, and then press \textbf{Escape}. Now you have two characters who ``belong''
    to Don Pedro.
  \item Use the same technique to add \textbf{Beatrice} and \textbf{Hero} underneath Leonato
    (they're his niece and daughter). Your diagram should appear as follows:

    \ogfigure{0.800}{res/og7_diagramming14.png}
\end{enumerate}

\section{Adding Labels to the Connecting Lines}

You can clarify the relationships between the characters by adding \textbf{Labels} to the
connecting lines. Do do this, you'll use the \textbf{Text} tool.

\ogfigure{0.800}{res/og7_diagrammingbasics_texttool.png}

\begin{enumerate}
  \item In the \textbf{Tool Palette}, double-click the \textbf{Text} tool---it's the one with a
    letter \textbf{A} on it---to make the tool ``sticky.''
  \item Place the mouse pointer over the line between Don Pedro and Claudio so that the line
    glows and then click it to make a label appear. Type \textbf{Boss} as the label's text to
    show that Claudio works for Don Pedro.

    \ogfigure{0.673}{res/og7_diagramming15.png}
  \item Add labels to the three remaining lines in the same way, as follows:

    \begin{itemize}
      \item Don Pedro is also Benedick's \textbf{Boss}
      \item Leonato is Hero's \textbf{Father} and Beatrice's \textbf{Uncle}
    \end{itemize}

    \ogfigure{0.800}{res/og7_diagramming16.png}
  \item Press \textbf{v} to turn off the Text tool and switch to the \textbf{Selection} tool.
\end{enumerate}

\section{Defining Relationships by Adding More Connections}

By now, you've already got a fine diagram, but sometimes you want to diagram relationships that
are not strictly tree-like. Let's manually add some more connections. Before you start
connecting things together manually, however, it's a good idea to turn off the automatic
layout feature.

\begin{enumerate}
  \item Press \textbf{Command--3} to open the \textbf{Canvas} inspectors.
  \item Option-click the \textbf{Diagram Layout and Style} inspector and turn off \textbf{Auto
    layout}:

    \ogfigure{0.800}{res/og7_diagramming17.png}

    With Auto layout disabled, the next part of the process is to further define the character's
    relationships by connecting their objects with the \textbf{Line} tool.
  \item Double-click the \textbf{Line} tool (or press \textbf{c} twice).
  \item Press \textbf{Command-1} to open to the \textbf{Object} inspectors.
  \item \textbf{Option-click} the \textbf{Line} inspector, and set the \textbf{Line type} to
    \textbf{Straight}:

    \ogfigure{0.800}{res/og7_diagramming18.png}

    \begin{ognote}
      With nothing selected on the canvas, if you choose the Shape, Line, Text, Pen, or Artboard
      tool and then make a style change in any of the Object inspectors, the changes you make
      become the \textit{new} default style for that tool. This applies not just to the current
      document, but to new documents or other document windows (or tabs) that you open in
      OmniGraffle.

      For example, in a new document, if you select the Shape tool and then choose the diamond
      shape in the Shape inspector, that sets the diamond as the default shape. If you had or
      were to open another OmniGraffle window, the shapes you draw with the Shape tool would
      also be diamonds.
    \end{ognote}
  \item Using the Line tool, make the following connections:

    \begin{itemize}
      \item Click \textbf{Don John} and drag a line to \textbf{Don Pedro}
      \item Click \textbf{Don Pedro} and drag a line to \textbf{Leonato}
      \item Click \textbf{Benedick} and drag a line to \textbf{Beatrice}
      \item Click \textbf{Claudio} and drag a line to \textbf{Hero}
    \end{itemize}
  \item Double-click the \textbf{Text} tool (or press \textbf{t} twice), and add labels to the
    lines as follows:

    \begin{itemize}
      \item Don John and Don Pedro are \textbf{Half-Brothers}
      \item Don Pedro and Leonato are \textbf{Friends}
      \item Benedick \textbf{Loves} Beatrice
      \item Claudio \textbf{Loves} Hero
    \end{itemize}

    Your diagram should appear as follows:

    \ogfigure{0.800}{res/og7_diagramming19.png}
  \item Press \textbf{v} to switch back to the \textbf{Selection} tool.
\end{enumerate}

\section{Using Orthogonal Lines to Denote Non-Linear Relationships}

Notice how the two ``Loves'' lines overlap each other and their respective objects; let's go
ahead and fix that.

\begin{enumerate}
  \item Using the \textbf{Selection} tool, \textbf{Command-click} to select the two ``Loves''
    lines.
  \item Press \textbf{Command-1} to view the \textbf{Object} inspectors.
  \item \textbf{Option-click} the \textbf{Line} inspector.
  \item Change the \textbf{Line type} from \textbf{Straight} to \textbf{Orthogonal}; however,
    nothing \textit{appears} to have changed\ldots{}but appearance isn't everything, right?

    \ogfigure{0.800}{res/og7_diagramming20.png}
  \item Click on the canvas to deselect the two lines, and then click on the line that extends
    from \textbf{Benedick} to \textbf{Beatrice}; notice the block of four squares on the path:

    \ogfigure{0.800}{res/og7_diagramming21.png}

    This block is a \textit{handle} which you can click and drag to reposition the orthogonal line
    so it routes around other objects in the diagram.
  \item Click and drag the \textbf{orthogonal line handle} to extend the line so that it
    highlights the relationship between Benedick and Beatrice:

    \ogfigure{0.800}{res/og7_diagramming22.png}
  \item Click on the canvas to deselect the line between Benedick and Beatrice, and then click
    to select the line between \textbf{Claudio} and \textbf{Hero}.
  \item Click and drag the \textbf{orthogonal line handle} to extend the line so that it
    highlights the relationship between Claudio and Hero; your diagram should look similar to
    the following:

    \ogfigure{0.800}{res/og7_diagramming23.png}
\end{enumerate}

\section{Adding Magnets to Objects}

By default, connecting lines are routed to the centerpoint of an object. For example, if you
select the line that connects Don John and Don Pedro and zoom in, you'll see a red and green
diamond at opposite ends of the lines:

\ogfigure{0.800}{res/og7_diagramming_magnets01.png}

A \textit{red diamond} denotes where the line started, while the \textit{green diamond} denotes
the line's terminus.

For most diagrams, routing the connecting lines to a centerpoint is effective. However, for
times when you want, or \textit{need}, the line to connect to a different place on the object,
you can add \textit{magnets} to an object and then re-route the connecting lines using the
\textbf{Selection} tool.

After adding the orthogonal lines that connect the \textbf{Love} relationships, you can achieve
a smoother line routing by adding magnets to the child objects in the diagram.

\begin{enumerate}
  \item Press and hold the \textbf{Command} key.
  \item Individually click to select \textbf{Benedick}, \textbf{Claudio}, \textbf{Beatrice}, and
    \textbf{Hero}:

    \ogfigure{0.800}{res/og7_diagramming_magnets02.png}
  \item Choose \textbf{Edit} $\blacktriangleright$ \textbf{Magnets} $\blacktriangleright$
    \textbf{Positions} $\blacktriangleright$ \textbf{8 Magnets}.
  \item To see where the magnets are positioned for the objects, choose \textbf{View}
    $\blacktriangleright$ \textbf{Magnets}:

    \ogfigure{0.800}{res/og7_diagramming_magnets03.png}

    The magnets are now visible as pink circles at each corner and midpoint along the child
    object's path, while the parent objects (Don John, Don Pedro, and Leonato) only have one
    magnet located at their centerpoint.
  \item Using the \textbf{Selection} tool (\textbf{v}), click to select the orthogonal line that
    connects \textbf{Claudio} to \textbf{Hero}, and reroute the line ends to the middle-bottom
    of each object:

    \ogfigure{0.800}{res/og7_diagramming_magnets04.png}

    \ogfigure{0.800}{res/og7_diagramming_magnets05.png}
  \item Drag the orthogonal line's handle downward to add more space between the objects and the
    line itself:

    \ogfigure{0.800}{res/og7_diagramming_magnets06.png}
  \item With the orthogonal line adjusted to your liking, you can turn off visibility of the
    magnets by again choosing \textbf{View} $\blacktriangleright$ \textbf{Magnets}.
\end{enumerate}

\section{Using the Line Inspector to Add a Line Hop}

When two lines pass over each other, such as the two orthogonal lines in this diagram, it
creates a slight possibility for someone to confuse which line goes where. The \textbf{Line}
inspector includes a pop-up menu from which you can specify a \textbf{Line Hop} for a selected
line.

\begin{enumerate}
  \item Click to select the orthogonal line between \textbf{Claudio} and \textbf{Hero}.
  \item \textbf{Option-click} the Line inspector.
  \item Click the \textbf{Line Hops} pop-up menu, and choose a style of your liking:

    \ogfigure{0.800}{res/og7_diagramming_linehops01.png}

    The pop-up menu includes options for specifying \textit{hops} that pass over or under another
    line to create tiny bridges or tunnels for the lines to pass through:

    \ogfigure{0.800}{res/og7_diagramming_linehops02.png}
\end{enumerate}

\section{Using the Object Inspectors to Add Styles to the Diagram}

You've already used a couple inspectors, but now you're going to use the inspectors in earnest.
There are 10 \textbf{Object} inspectors, and each one controls a certain aspect of the
selected objects' appearance. The first inspector you'll use to style the shapes in the
diagram is the \textbf{Fill} inspector.

\begin{enumerate}
  \item Select \textbf{Don John} on the canvas.
  \item Press \textbf{Command-1} to view the \textbf{Object} inspectors.
  \item \textbf{Option-click} the \textbf{Fill} inspector, which reflects that the \textbf{Don
    John} object has a plain white fill color:

    \ogfigure{0.800}{res/og7_diagramming24.png}
  \item Click the \textbf{Fill Type} pop-up menu, and choose \textbf{Linear Gradient} from the
    list of available fill types.
  \item Beneath the pop-up menu, you will see two color menus for defining the gradient's
    colors:

    \begin{itemize}
      \item Click the \textbf{Beginning gradient color} menu on the left and choose a color.
      \item Click the \textbf{Ending gradient color} menu on the right and choose another,
        complementary, color.
    \end{itemize}

    \ogfigure{0.800}{res/og7_diagramming25.png}
  \item \textbf{Option-click} the \textbf{Stroke} inspector.
  \item Change the thickness of the stroke to \textbf{2 pt}.
  \item \textbf{Option-click} the \textbf{Shapes} inspector.
  \item Change the \textbf{Corner radius} to \textbf{5 pt}:

    \ogfigure{0.800}{res/og7_diagramming27.png}

    This makes for a nice, heavy, rounded stroke around the object. You can zoom in so you can see
    the progress you've made!
  \item \textbf{Option-click} the \textbf{Font} inspector.
  \item Choose a different font, then click the color well next to the Font button and choose a
    color to contrast with the fill. Now you should have one quite charming shape object.

    \ogfigure{0.800}{res/og7_diagramming28.png}
\end{enumerate}

\section{Using Style Swatches to Rapidly Apply Styles}

With Don John still selected, look at the bottom of the Inspector bar for the \textbf{Style
Tray}:

\ogfigure{0.746}{res/og7_diagramming_styletray.png}

The Style Swatch (or \textit{chit}) on the left combines all of the style properties of the
selected object, while the grouping of Style Swatches on the right lists them individually by
inspector.

\begin{enumerate}
  \item Click and drag the \textbf{All Styles} Swatch and drop it onto the \textbf{Don Pedro}
    shape.

    \ogfigure{0.800}{res/og7_diagramming29.png}

    \ogfigure{0.800}{res/og7_diagramming30.png}
  \item Do the same thing to \textbf{Leonato} so that all of the elder generation characters
    have a consistent appearance.

    \ogfigure{0.800}{res/og7_diagramming31.png}
\end{enumerate}

\section{Selecting and Styling the Child Objects}

Let's style the remaining objects. A quick way to select lots of shapes at once is with the
\textbf{Selection} tab of the Sidebar (\textbf{Option-Command-5}). The \textbf{Selection} tab
lists all of the different kinds of objects on the canvas---lines, labels, shapes, and so
on---based on their styles.

\begin{enumerate}
  \item In the \textbf{Selection} tab, click the item representing the four plain white
    rectangles with black strokes. This selects all four objects at once.

    \ogfigure{0.800}{res/og7_diagramming32.png}
  \item Use the \textbf{Shape} inspector to change the shape of the four selected objects from
    Rectangles to Circles (or, in this case, Ellipses).

    \ogfigure{0.800}{res/og7_diagramming33.png}
  \item Use the \textbf{Fill}, \textbf{Stroke}, and \textbf{Font} inspectors to simultaneously
    style all four shapes in in a way that differentiates the three rectangular shapes. Here
    we've chosen a different color, a slightly thicker stroke, and a larger font.

    \ogfigure{0.800}{res/og7_diagramming34.png}
\end{enumerate}

\section{Adding Layers to Your Project}

Layers contain different sets of objects on the same canvas. We'll place a background image on a
new layer to keep it separate from the character diagram.

\begin{enumerate}
  \item Before adding a new layer to your project, choose \textbf{View} $\blacktriangleright$
    \textbf{Show Contents} $\blacktriangleright$ \textbf{Layers} (or press
    \textbf{Option-Command-2}) to switch to the Layers tab in the Sidebar.
  \item In the toolbar, click \textbf{New Layer} to add another one to the canvas. A new layer
    is born. You can name it if you like; something like ``Background'' should suffice.

    \ogfigure{0.673}{res/og7_diagramming35.png}
  \item As you learned in The Left Sidebar, the order of layers is important; objects on higher
    layers appear in front of stuff on lower layers. Since you want the Background layer and its
    image to appear behind the diagram, drag the Diagram layer above the Background layer.

    \ogfigure{0.800}{res/og7_diagramming36.png}
  \item Let's concentrate on the Background layer. The layer you're working with right now has a
    pencil icon to the left of its preview in the Sidebar. If the Background layer doesn't have
    the pencil icon, click just to the left of the layer preview to put it there.

    \ogfigure{0.673}{res/og7_diagramming37.png}
  \item For your background, rather than using the \textbf{Background Image} inspector, which
    would place a background image on the entire canvas, you'll want to drag an image onto the
    layer. A new shape containing the image is created on the Background layer.

    \ogfigure{0.800}{res/og7_diagramming38.png}
  \item Use the controls in the \textbf{Image} inspector to change the size and scaling or the
    opacity of the image to suit your liking.

    \ogfigure{0.800}{res/og7_diagramming39.png}

    \ogfigure{0.800}{res/og7_diagramming40.png}
\end{enumerate}

Now you have an attractive, informative diagram.

\ogfigure{0.800}{res/og7_diagramming41.png}

As you can see, there's still some work to be done on the connecting lines and their labels, so
let's move on to those next.

\section{Styling Lines and Labels}

So far, we've focused on styling the object blocks and giving the diagram a nice background. But
if you look at the previous image, you can see that there's still some work to do,
particularly with the lines and labels.

Here we want the lines to help people visualize the relationships between the characters, and to
do more than just show a simple connection. So let's address the lines first, and then you'll
see how to remove the fill from a label.

\begin{enumerate}
  \item Choose \textbf{Inspectors} $\blacktriangleright$ \textbf{Object} (\textbf{Command--1})
    to view the Object inspectors on the right.
  \item \textbf{Option-click} the \textbf{Stroke} inspector.
  \item Select the line between Don John and Don Pedro.
  \item In the \textbf{Stroke} inspector, set the stroke width to \textbf{1.5 pt}, and change
    the line's color to a nice shade of blue.
  \item With the line still selected, \textbf{Option-click} on the \textbf{Line} inspector; the
    Stroke inspector collapses while the Line inspector opens.
  \item Give the line arrows at both the left and right sides. Your diagram and inspectors
    should look something like this:

    \ogfigure{0.800}{res/og7_diagramming42.png}
  \item Use a similar approach for the other lines, as follows:

    \begin{itemize}
      \item Create a double-arrowed line between Don Pedro and Leonato
      \item Create single-arrowed lines that point up from Claudio and Benedick to Don Pedro
      \item Create single-arrowed lines that point up from Beatrice and Hero to Leonato
      \item Create double-arrowed lines between Benedick and Beatrice, and between Claudio and
        Hero
    \end{itemize}

    When you've finished, your diagram should look something like the following:

    \ogfigure{0.800}{res/og7_diagramming43.png}
\end{enumerate}

Now let's deal with the labels on the lines. As you can see, they have a white background that
obscures the lines and blots out the background image. Let's fix that!

\begin{enumerate}
  \item To make this process go faster, let's go back to our old friend, the \textbf{Selection}
    tab (\textbf{Option-Command--5}). The Selection tab, if you recall from Selecting and
    Styling the Child Objects, lets you quickly select similar items in your diagram. In this
    case, you'll want to select all eight of the text labels in the Selection tab.

    \ogfigure{0.800}{res/og7_diagramming44.png}
  \item Choose \textbf{Inspectors} $\blacktriangleright$ \textbf{Object} (\textbf{Command--1}),
    \textbf{Option-click} on the \textbf{Fill} inspector, and then set the \textbf{Fill} for the
    labels to \textbf{No Fill}. This removes the white box behind the labels so the lines and
    the background image show through.

    \ogfigure{0.800}{res/og7_diagramming45.png}

    The only problem is that---with some of the labels---the text interferes with the lines. Let's
    fix that.
  \item In the diagram, \textbf{Command-click} the \textbf{Love} labels to select both of those.
  \item \textbf{Option-click} the \textbf{Geometry} inspector to change their positioning
    related to the line.

    \ogfigure{0.800}{res/og7_diagramming46.png}
  \item Select the \textbf{Half-Brothers} label and tap the \textbf{Up Arrow} key on your
    keyboard to nudge the label above the line.
  \item Select the \textbf{Friends} label and nudge that above the line using the \textbf{Up
    Arrow} key on your keyboard.
  \item Use \textbf{Command-click} to select the two \textbf{Boss}, \textbf{Uncle}, and
    \textbf{Father} labels.
  \item In the \textbf{Geometry} object inspector, change their orientation relative to the line
    from \textbf{Horizontal} to \textbf{Parallel}.

    \ogfigure{0.800}{res/og7_diagramming47.png}
  \item Finally, select each of those four labels individually to move them off the lines. Your
    diagram should look similar to the following:

    \ogfigure{0.800}{res/og7_diagramming48.png}
\end{enumerate}

\section{Adding a Title to the Diagram}

Now it's time for you to take a little bit of a breather. You've been going at this pretty hard
now, so we'll make this one simple. All you need to do is add a title to the diagram.

\begin{enumerate}
  \item Choose \textbf{View} $\blacktriangleright$ \textbf{Show Contents} $\blacktriangleright$
    \textbf{Layers} (or press \textbf{Option-Command-2}) to switch to the Layers tab in the
    Sidebar.
  \item In the toolbar, click \textbf{New Layer} to add another layer to the project. Name this
    layer \textbf{Title}:

    \ogfigure{0.800}{res/og7_diagramming49.png}
  \item Select the \textbf{Text} tool, click on the Canvas in the open area to the left of
    \textbf{Don John}, and enter \textit{Much Ado About Nothing}.
  \item Use the \textbf{Font} and \textbf{Text Position} inspectors to style the title however
    you'd like.

    \ogfigure{0.800}{res/og7_diagramming50.png}
\end{enumerate}

\section{Where to Go From Here}

When you've finished, here's what the \textit{Much Ado About Nothing} diagram looks like:

\ogfigure{0.800}{res/og7_diagramming51.png}

Take some time now to experiment with OmniGraffle a little more. The more you use OmniGraffle,
the more comfortable you'll become with all of its abilities. Don't worry about ``messing''
anything up; just remember that you can always press \textbf{Command-Z} to undo any recent
changes you've made.
%%%%%%%%%%%%%%%%%%%%%%%%%%%%%%%%%%%%%%%%%%%%%%%%%%%%%%%%%%%%%%%%%%%%%%%%%%%%%%%%%%%%%%%%%%%%%%%%%%%
\end{document}
